\documentclass[11pt]{article}
\usepackage[a4paper,margin=25mm]{geometry}
\usepackage[T1]{fontenc}
\usepackage{lmodern,microtype}
\usepackage{amsmath,amssymb,amsthm,mathtools}
\usepackage{enumitem,booktabs,xurl}
\usepackage[hidelinks]{hyperref}
\newcommand{\NE}{\operatorname{NE}}
\newcommand{\BR}{\operatorname{BR}}
\newcommand{\Stab}{\operatorname{Stab}_0}
\newcommand{\supp}{\operatorname{supp}}
\newcommand{\rank}{\operatorname{rank}}
\newcommand{\im}{\operatorname{im}}
\newcommand{\relint}{\operatorname{relint}}
\newcommand{\Q}{\mathbb Q}
\newcommand{\R}{\mathbb R}
\newcommand{\one}{\mathbf 1}
\newtheorem{theorem}{Theorem}[section]
\newtheorem{lemma}[theorem]{Lemma}
\newtheorem{proposition}[theorem]{Proposition}
\newtheorem{corollary}[theorem]{Corollary}
\theoremstyle{definition}
\newtheorem{definition}[theorem]{Definition}
\newtheorem{example}[theorem]{Example}
\theoremstyle{remark}
\newtheorem{remark}[theorem]{Remark}
\setlist[enumerate]{itemsep=3pt,topsep=4pt}
\setlength{\emergencystretch}{2em}
\hypersetup{pdftitle={Two-player partial results for axiomatic containment of Mertens-stable sets},pdfauthor={LLM-assisted research draft}}
\title{Two-player partial results for axiomatic containment\\
 of Mertens-stable sets without genericity\\[4pt]
 \large We-Solve problem Econ-20 (Economics rank Econ-20)}
\author{LLM-assisted research contribution}
\date{19 October 2026}
\begin{document}
\maketitle

\begin{abstract}
We study the restriction to finite bimatrix target games of the question
whether Govindan--Wilson's global axioms $B,D^*,I^*$ compel every
selected nonempty compact connected Nash-equilibrium set to contain a
Mertens characteristic-zero stable set. We prove the containment
conclusion for all $2\times2$ games, give a rank criterion for the
stability of quasi-strict bimatrix equilibrium singletons, establish a
finite cross-game exclusion certificate and use it to prove stability
of selected sets consisting entirely of quasi-strict equilibria. We
settle a degenerate $2\times3$ example at the set level; we prove the
conclusion for all $2\times n$ games with a strictly dominant column;
and a calibrated dominated-action construction forces the exact
selection of an entire equilibrium face when an explicit balancing and
spanning condition holds. The latter construction includes equilibrium
simplices of unbounded dimension. These are \emph{partial results}:
no proof or counterexample for all finite two-player games is claimed.
\end{abstract}

\section*{Submission identification and claim}\label{sec:metadata}
\noindent\textbf{Public submission handle:} \texttt{Econ-20}.\\
\textbf{Catalogue ranking label:} Econ-20 (not a permanent ID).\\
\textbf{Research-record ID:} \texttt{prb\_34ecf5fa262f4a85a0fbe4af9837c46b}.\\
\textbf{Earlier draft label:} \texttt{WS-107793} (not verified as a
public alias of the current We-Solve record).\\
\textbf{Claim in contribution form:} \emph{Partial result}.\\
\textbf{Scope:} All $2\times2$ target games, selected sets consisting
entirely of quasi-strict bimatrix equilibria, the stated dominant-column
and balanced--spanning subclasses, and the stated exact exclusion
certificates. The tested game is restricted; the axioms are always
assumed \emph{globally} on all finite games. Arbitrary finite bimatrix
targets remain unresolved.

\medskip
\noindent\textbf{Source and attribution.} This manuscript consolidates
three LLM-assisted research drafts supplied by the submitter. The
mathematical arguments are presented in full rather than by references
to those drafts. The definitions, all-game nonemptiness and axiom
satisfaction of the Mertens stability solution concept, and the generic
characterization are credited to Govindan--Wilson~\cite{gw}; none of
the results below claims literature priority or independent referee
verification. The submitter should report \emph{all} actual LLMs and
humans involved, with real thinking settings or ``unknown'', separately
in the We-Solve issue form; this source text does not infer settings.

\section{Exact two-player target and source inputs}\label{sec:scope}
Let $\Phi$ be a rule on \emph{all finite normal-form games}, assigning to each
$G$ a nonempty family of nonempty closed connected subsets of $\NE(G)$.
Assume that $\Phi$ satisfies the exact axioms $B,D^*,I^*$ in the pinned
Govindan--Wilson paper \cite[Sections 3--5]{gw}. The two-player question is
\begin{equation}\label{eq:target}
 \forall\,G=(A,B)\text{ finite bimatrix game}\quad
 \forall K\in\Phi(G)\quad
 \exists S\in\Stab(G)\quad S\subseteq K.
\end{equation}
The sizes $m\times n$ are arbitrary. Only the target game is restricted to two
players: the domain of $\Phi$ and the comparisons allowed by the axioms are
\emph{not} restricted to two players. The source proves the corresponding
containment theorem for generic extensive-form games and discusses the
nongeneric obstruction \cite[Theorem 6.1 and Section 10]{gw}.

For the positive $2\times2$ theorem we need only this consequence of axiom $B$:
\begin{equation}\label{eq:Bsim}
 K\in\Phi(G)\quad\Longrightarrow\quad
 K\text{ contains a normal-form trembling-hand perfect equilibrium.}
\end{equation}
Indeed, represent the normal-form game by one simultaneous decision per player,
with the second player not observing the first player's action. Each has just
one information set. Under the source's Definition 2.1, quasi-perfection in
this representation is the usual normal-form perfection condition. There are
no subsequent own decisions. Keeping the terminal action profiles distinct
makes realization equivalence identify the same mixed strategies.

\begin{definition}[Perfection used in this note]
A profile $(p^*,q^*)$ is perfect if there are completely mixed profiles
$(p^k,q^k)\to(p^*,q^*)$ such that $p^*$ is a best reply to $q^k$ and $q^*$ is a
best reply to $p^k$ for every $k$. In finite games this is the usual
best-reply characterization of normal-form trembling-hand perfection.
\end{definition}


\noindent\textbf{Independent known inputs.} Govindan--Wilson show that
Mertens's characteristic-zero stability assigns a nonempty solution
family to each finite game and satisfies $B,D^*,I^*$; their converse
containment theorem is established for generic games only
\cite[Theorem 6.1 and Appendices A--B]{gw}. These facts are used below
only to identify stable-set families \emph{after} deriving an exact
selection statement, never to assume nongeneric containment.

\section{Strategy perturbations and two topological lemmas}
For a finite game with mixed space $\Sigma=\prod_i\Delta(A_i)$, put
\[
 P_\epsilon=\{\delta\tau:0\le\delta\le\epsilon,\ \tau\in\Sigma\}.
\]
If $\eta\in P_1$, the common mass $\delta(\eta)=\sum_a\eta_i(a)$ does not depend
on the player. The perturbed feasible set is
\[
 \Sigma_i(\eta)=\{x_i\in\Delta(A_i):x_i(a)\ge\eta_i(a)\text{ for every }a\}.
\]
Let $E_G$ be the graph of Nash equilibria of these constrained games over
$P_1$, with projection $\pi$. Boundaries of $P_\epsilon$ are relative to its
linear span, which has dimension
\[
 d=1+\sum_i(|A_i|-1).
\]
The pair $(P_\epsilon,\partial P_\epsilon)$ is homeomorphic to
$(D^d,S^{d-1})$. For $X\subseteq E_G$, write
\[
 X_\epsilon=X\cap\pi^{-1}(P_\epsilon),\qquad
 \partial X_\epsilon=X\cap\pi^{-1}(\partial P_\epsilon).
\]
The notation $\partial X_\epsilon$ means this \emph{inverse image}, not an
intrinsic boundary chosen for $X$.

We use the standard closed-witness formulation of Mertens $0$-stability:
a closed $X\subseteq E_G$ with $X_0=\{0\}\times S$, nonzero pullback
\[
 \pi^*:\check H^d(P_\epsilon,\partial P_\epsilon;\Q)
       \longrightarrow\check H^d(X_\epsilon,\partial X_\epsilon;\Q)
\]
for a positive $\epsilon$, and the source's connexity condition. Connexity
requires that, for every neighborhood $V$ of $X_0$ in $X$, some component of
$V\setminus\partial X_1$ have closure a neighborhood of $X_0$ in $X$.
These are the strategy-perturbation and rational-cohomology conventions of
\cite[Definition 2.4]{gw}, with the closed-witness convention made explicit;
see also the definition framework of Govindan--Mertens \cite{gm}.

\begin{lemma}[Continuous equilibrium section]\label{lem:section}
Suppose $z\in\NE(G)$ and, for some $0<\rho<1$, a continuous map
$s:P_\rho\to\Sigma$ satisfies
\[
 s(0)=z,\qquad s(\eta)\in\NE(G(\eta))\quad(\eta\in P_\rho).
\]
Then $\{z\}\in\Stab(G)$.
\end{lemma}
\begin{proof}
Take $X=\{(\eta,s(\eta)): \eta\in P_\rho\}$. This is compact, is contained in
$E_G$, and has zero fiber $\{(0,z)\}$. For every $0<\epsilon\le\rho$, projection
is a homeomorphism of pairs
\[
 (X_\epsilon,\partial X_\epsilon)\longrightarrow
 (P_\epsilon,\partial P_\epsilon).
\]
The induced degree-$d$ map is therefore an isomorphism between copies of
$\Q$, and in particular is nonzero.

Let $V$ be a neighborhood of $(0,z)$ in $X$. There is $0<\kappa<\rho$ such
that the graph over $P_\kappa$ lies in $V$. The graph over
\[
 U_\kappa=\{\eta:0<\delta(\eta)<\kappa,\ \eta_i(a)>0\ \forall i,a\}
\]
is connected: its base is homeomorphic to
$(0,\kappa)\times\relint\Sigma$. It lies in $V\setminus\partial X_1$.
Strictly positive directions are dense in $P_\kappa$, so its closure contains
the graph over $P_{\kappa/2}$. Consequently, the component of
$V\setminus\partial X_1$ containing that graph has closure containing a
neighborhood of $(0,z)$ in $X$. This verifies connexity separately from
cohomological essentiality.
\end{proof}

\begin{lemma}[Necessary local robustness]\label{lem:necessary}
If $S\in\Stab(G)$, then for every open neighborhood $W$ of $S$ there is
$\rho>0$ such that every $\eta\in P_\rho$ with $0<\delta(\eta)$ and all
coordinates strictly positive has a constrained equilibrium in $W$.
\end{lemma}
\begin{proof}
Let $X$ be a closed stability witness and choose an $\epsilon>0$ for which
the cohomological pullback is nonzero. Its projection contains every interior
point $a$ of $P_\epsilon$. Otherwise it factors through
\[
 (P_\epsilon\setminus\{a\},\partial P_\epsilon),
\]
whose relative degree-$d$ cohomology is zero: radial projection away from $a$
retracts the punctured polytope onto its boundary. The induced pullback would
then be zero, a contradiction.

Next, all witness equilibria over sufficiently small common masses lie in
$W$. If not, there are $(\eta^k,z^k)\in X$ with $\eta^k\to0$ and
$z^k\notin W$. Compactness supplies a subsequence converging to $(0,z)\in X$
with $z\notin W$. But the zero fiber is $S\subset W$. This is impossible.
Combining this observation with surjectivity onto the interior of
$P_\epsilon$ proves the assertion. Essentiality at one $\epsilon$ suffices;
no unproved passage of essentiality to smaller witnesses is used.
\end{proof}


\section{A sufficient stability witness for a connected set}
Let $\Sigma=\prod_i\Delta(A_i)$ and
\[
 P_\rho=\{\delta\tau:0\le\delta\le\rho,\ \tau\in\Sigma\},\qquad
 \Sigma_i(\eta)=\{x_i\in\Delta(A_i):x_i(a)\ge\eta_i(a)\}.
\]
The common mass is $\delta(\eta)=\sum_a\eta_i(a)$, independent of $i$.
Write $E_G$ for the equilibrium graph of the constrained games and
$\pi(\eta,x)=\eta$. The perturbation polytope has dimension
$d=1+\sum_i(|A_i|-1)$; its boundary is relative to its linear span.
For a witness $X\subseteq E_G$, the symbol $\partial X_\epsilon$ means
$X\cap\pi^{-1}(\partial P_\epsilon)$, not an independently chosen boundary.

\begin{lemma}[A continuous family of sections]\label{lem:family}
Let $K\subseteq\NE(G)$ be nonempty, compact and connected. Suppose that,
for some $0<\rho<1$, a jointly continuous map
\[
 s:K\times P_\rho\longrightarrow\Sigma
\]
satisfies $s(z,0)=z$ and $s(z,\eta)\in\NE(G(\eta))$ for every
$(z,\eta)$. Then $K\in\Stab(G)$. Moreover, every $z\in K$ has a
stable singleton.
\end{lemma}
\begin{proof}
Let $X=\{(\eta,s(z,\eta)):z\in K,\eta\in P_\rho\}$.
It is compact, lies in $E_G$, and has zero fiber $\{0\}\times K$.
Fix $z_0\in K$. For $0<\epsilon\le\rho$, the map
\[
 j:(P_\epsilon,\partial P_\epsilon)\longrightarrow
       (X_\epsilon,\partial X_\epsilon),
 \qquad j(\eta)=(\eta,s(z_0,\eta))
\]
satisfies $\pi\circ j=\operatorname{id}$. Hence $j^*\pi^*$ is the
identity on $\check H^d(P_\epsilon,\partial P_\epsilon;\Q)\cong\Q$,
so $\pi^*$ is nonzero.

To check connexity separately, let $V$ be any neighborhood of $X_0$ in
$X$. Compactness implies that $X_\kappa\subseteq V$ for some
$0<\kappa<\rho$: otherwise a sequence outside $V$ with masses tending
to zero has a limit in $X_0$, contradicting the neighborhood property.
The set
\[
 U_\kappa=\{\eta:0<\delta(\eta)<\kappa,\ \eta_i(a)>0\text{ for all }i,a\}
\]
is connected, being homeomorphic to
$(0,\kappa)\times\relint\Sigma$. The image of $K\times U_\kappa$
is connected and lies in $V\setminus\partial X_1$. Its closure contains
$X_{\kappa/2}$, by density of strictly positive perturbations and continuity
of $s$. A component of $V\setminus\partial X_1$ containing this image
therefore has closure containing a neighborhood of $X_0$. This is the
required connexity. Restricting $s$ to any singleton parameter $z$ proves
the last assertion by the same argument.
\end{proof}

\section{\texorpdfstring{Complete containment theorem for all $2\times2$ games}{Complete containment theorem for all 2 by 2 games}}
Represent each player's mixed strategy by the probability of action $1$:
$x,y\in[0,1]$. The two payoff differences are affine functions
\[
 f(y)=u_1(1,y)-u_1(0,y),\qquad
 g(x)=u_2(x,1)-u_2(x,0).
\]

\begin{theorem}[Every perfect equilibrium of a $2\times2$ game is singleton-stable]
\label{thm:2by2}
For every finite $2\times2$ game, every normal-form perfect equilibrium
$z^*=(x^*,y^*)$ admits a local continuous equilibrium section of its
lower-bound perturbations. Hence
\[
 \{z^*\}\in\Stab(G).
\]
No genericity, isolation, strictness, or nondegeneracy assumption is made.
\end{theorem}
\begin{proof}
For the row player define $\ell_1=\eta_{1,1}$ and
$u_1=1-\eta_{1,0}$, and define $\ell_2,u_2$ similarly. Thus the perturbed
strategy intervals are $[\ell_1,u_1]$ and $[\ell_2,u_2]$. Set
\begin{equation}\label{eq:clamp}
 x^\eta=\min\{u_1,\max\{\ell_1,x^*\}\},\qquad
 y^\eta=\min\{u_2,\max\{\ell_2,y^*\}\}.
\end{equation}
This map is continuous, feasible, and equals $z^*$ at zero.
For sufficiently small common mass, any interior target coordinate is
unchanged, while a target $0$ is sent to its lower bound and a target $1$
is sent to its upper bound.

We prove that $x^\eta$ is a constrained best reply to $y^\eta$.
The column-player argument uses $g$ and interchanges the roles of $x,y$.

\smallskip
\noindent\emph{Case 1: $0<y^*<1$.}
Choose the mass small enough that $y^\eta=y^*$. The Nash condition gives
$f(y^*)\le0$ if $x^*=0$, $f(y^*)\ge0$ if $x^*=1$, and $f(y^*)=0$ if
$0<x^*<1$. In the first two cases the clamped $x^\eta$ is the appropriate
feasible endpoint. In the last case every feasible $x$ is optimal.

\smallskip
\noindent\emph{Case 2: $y^*\in\{0,1\}$ and $0<x^*<1$.}
Perfection supplies $y^k\in(0,1)$ tending to $y^*$, with $x^*$ a best reply
for every $k$. Since $x^*$ uses both actions, $f(y^k)=0$ for every $k$.
A nonzero affine function cannot vanish at such a sequence approaching a
boundary point. Therefore $f$ is identically zero, and every feasible
$x^\eta$ is a best reply.

\smallskip
\noindent\emph{Case 3: $x^*,y^*\in\{0,1\}$.}
If $x^*=1$, perfection gives $f(y^k)\ge0$ for an interior sequence approaching
$y^*$. An affine function with this property is nonnegative on a whole
one-sided neighborhood of $y^*$. To check this explicitly, if $f(y^*)>0$
use continuity; if $f(y^*)=0$, the sign just inside the interval is fixed by
its slope; and if $f(y^*)<0$, the sequence is impossible. Thus $f(y^\eta)\ge0$
for all sufficiently small bounds, and $x^\eta=u_1$ is optimal. When $x^*=0$
the same reasoning uses the weak inequality $f\le0$, and $x^\eta=\ell_1$
is optimal.

There are only two players and finitely many neighborhood requirements.
Taking their minimum yields one positive $\rho$ valid for every perturbation
in $P_\rho$, including directions with some zero coordinates. Both players
are constrained best replies. Lemma~\ref{lem:section} proves singleton stability.
\end{proof}

\begin{corollary}[Econ-20 for every $2\times2$ target]\label{cor:axioms2by2}
Let $\Phi$ satisfy axiom $B$ globally, and let $G$ be any $2\times2$ game.
Every selected set $K\in\Phi(G)$ contains a Mertens $0$-stable singleton.
In particular, \eqref{eq:target} holds for all $2\times2$ targets under the
original assumptions $B,D^*,I^*$.
\end{corollary}
\begin{proof}
By \eqref{eq:Bsim}, $K$ contains a perfect equilibrium $z^*$. Apply
Theorem~\ref{thm:2by2} and take $S=\{z^*\}$.
\end{proof}

\begin{remark}
The containment argument does not need connectedness of $K$, $D^*$, or $I^*$.
It uses only the simultaneous-representation consequence of $B$ and
nonemptiness of the selected family. A theorem for $2\times2$ games is not a
theorem for arbitrary two-player games. It is also not a theorem for all
extensive-form games having two actions per information set: such games can
have many pure contingent plans.
\end{remark}


\section{An exact singleton criterion at quasi-strict bimatrix equilibria}
Let the two payoff matrices be $A,B\in\R^{m\times n}$. The row player uses
$p\in\Delta_m$ and the column player $q\in\Delta_n$. Let $(p^*,q^*)$ be a
Nash equilibrium with supports $I=\supp(p^*)$ and $J=\supp(q^*)$.

\begin{definition}[Quasi-strictness]
The equilibrium is quasi-strict if the row player's pure best replies to
$q^*$ are exactly $I$, and the column player's pure best replies to $p^*$
are exactly $J$. Thus every unused action is strictly inferior against the
target opponent. This is not the definition of quasi-perfect equilibrium.
\end{definition}

Choose $i_0\in I$ and $j_0\in J$. Define matrices of supported payoff
differences by
\[
 C_{i,:}=A_{i,:}-A_{i_0,:}\quad(i\in I\setminus\{i_0\}),\qquad
 D_{j,:}=(B_{:,j}-B_{:,j_0})^{\mathsf T}\quad(j\in J\setminus\{j_0\}).
\]
An empty collection of difference rows is allowed. Put
\begin{equation}\label{eq:rankmatrices}
\begin{aligned}
 M_A&=\begin{pmatrix}C_{:,J}\\\one_J^{\mathsf T}\end{pmatrix},&
 N_A&=\begin{pmatrix}C_{:,J^c}\\\one_{J^c}^{\mathsf T}\end{pmatrix},\\
 M_B&=\begin{pmatrix}D_{:,I}\\\one_I^{\mathsf T}\end{pmatrix},&
 N_B&=\begin{pmatrix}D_{:,I^c}\\\one_{I^c}^{\mathsf T}\end{pmatrix}.
\end{aligned}
\end{equation}
Here $M_A$ is $|I|\times|J|$ and $M_B$ is $|J|\times|I|$.
An off-support matrix has zero columns when its complement is empty.

\begin{theorem}[Rank characterization at a quasi-strict equilibrium]
\label{thm:rank}
At a quasi-strict bimatrix equilibrium $z^*=(p^*,q^*)$, the following are
equivalent:
\begin{enumerate}[label=(\roman*)]
\item $\{z^*\}\in\Stab(G)$.
\item Every neighborhood $W$ of $z^*$ contains an equilibrium of every
sufficiently small strictly positive common-mass lower-bound perturbation.
\item
\begin{equation}\label{eq:rankcondition}
 \im N_A\subseteq\im M_A,\qquad \im N_B\subseteq\im M_B.
\end{equation}
Equivalently, $\rank M_A=\rank[M_A\ N_A]$ and
$\rank M_B=\rank[M_B\ N_B]$.
\item The constrained-equilibrium correspondence has a local affine section
through $z^*$, defined for all sufficiently small nonnegative lower bounds.
\end{enumerate}
\end{theorem}
\begin{proof}
The implication (i)$\Rightarrow$(ii) is Lemma~\ref{lem:necessary}.
We first prove (iii)$\Rightarrow$(iv). Choose matrices $R_A,R_B$ with
\[
 M_A R_A=N_A,\qquad M_B R_B=N_B.
\]
For row lower bounds $\alpha$ and column lower bounds $\beta$, define
\begin{equation}\label{eq:affinesection}
\begin{aligned}
 q_{J^c}&=\beta_{J^c},&q_J&=q_J^*-R_A\beta_{J^c},\\
 p_{I^c}&=\alpha_{I^c},&p_I&=p_I^*-R_B\alpha_{I^c}.
\end{aligned}
\end{equation}
The last row of each matrix identity imposes total mass one. The other
rows give $Cq=0$ and $Dp=0$, so all originally supported actions remain
mutually indifferent for their respective players.

Let $\mu$ be the smallest positive coordinate of $p_I^*$ and $q_J^*$.
It is positive because the supports are finite and nonempty. The changes in
supported probabilities in \eqref{eq:affinesection} are bounded by a constant
times $\delta(\alpha,\beta)$. Choosing that mass sufficiently small makes every
supported coordinate exceed $\mu/2$ and all lower bounds less than $\mu/2$.
Thus feasibility holds. Quasi-strictness and continuity preserve every strict
payoff gap for unused actions. A player's constrained best reply must assign
the lower bound to each strictly inferior action and may distribute the
remaining probability among the tied best actions. That is exactly what
\eqref{eq:affinesection} does. This proves (iv). Its continuity and
Lemma~\ref{lem:section} give (iv)$\Rightarrow$(i).

It remains to prove (ii)$\Rightarrow$(iii). Choose a fixed neighborhood $W$
of $z^*$ in which all coordinates on $I,J$ exceed $\mu/2$, and each unused
action remains strictly inferior to each originally supported action.
Such a neighborhood exists by quasi-strictness and continuity.
For bounds with common mass below $\mu/2$, any constrained equilibrium in
$W$ must have
\[
 p_{I^c}=\alpha_{I^c},\quad q_{J^c}=\beta_{J^c},\quad Cq=0,\quad Dp=0.
\]
The supported coordinates are strictly above their lower bounds, so they
must tie at the maximal payoff; inferior coordinates must be at their
bounds. Subtracting the target equations gives the necessary systems
\begin{equation}\label{eq:necessarysystems}
 M_A(q_J-q_J^*)=-N_A\beta_{J^c},\qquad
 M_B(p_I-p_I^*)=-N_B\alpha_{I^c}.
\end{equation}

Suppose, for instance, $\im N_A\not\subseteq\im M_A$. Finite-dimensional
linear algebra gives a row vector $w$ with $wM_A=0$ but $wN_A\ne0$.
There is a vector $b$ with all coordinates strictly positive and
$wN_A b\ne0$: a nonzero linear form cannot vanish on the open positive
orthant. Scale $b$ so $\sum b<1$, extend it on $J$ by strictly positive
coordinates of total mass $1-\sum b$, and call the resulting full
probability vector $\bar\beta$. Choose any completely mixed
$\bar\alpha\in\Delta_m$.
For every $\delta>0$ set $\alpha=\delta\bar\alpha$ and
$\beta=\delta\bar\beta$. These are strictly positive perturbations with
common mass $\delta$, but multiplying the first equation of
\eqref{eq:necessarysystems} by $w$ gives
\[
 0=-\delta wN_A b\ne0.
\]
No constrained equilibrium lies in $W$, however small $\delta$ is. This
contradicts (ii). The second inclusion is proved by the same argument.
\end{proof}

\begin{corollary}[Isolated quasi-strict equilibria]\label{cor:isolated}
Every isolated quasi-strict equilibrium of a finite bimatrix game has a
Mertens $0$-stable singleton.
\end{corollary}
\begin{proof}
Suppose $M_A$ has a nonzero kernel vector $v$. The profile with $p=p^*$,
$q_{J^c}=0$, and $q_J=q_J^*+tv$ is a Nash equilibrium for every sufficiently
small $|t|$: supported probabilities stay positive; the last matrix row
preserves their sum; the difference rows preserve row-player indifference;
and the unused-action strict inequalities persist. Column best replies are
unchanged since $p=p^*$. This contradicts isolation. Thus $M_A$ has full
column rank, so $|J|\le|I|$. The same argument for $M_B$ gives
$|I|\le|J|$. Consequently the supports have equal size and both $M_A,M_B$
are square nonsingular matrices. Their column spaces are their full target
spaces, so \eqref{eq:rankcondition} holds. Apply Theorem~\ref{thm:rank}.
\end{proof}

\begin{remark}[What the local criterion does not say]
Quasi-strictness alone does not imply the rank conditions, as the next
example shows. Nor does the criterion imply that every axiomatically selected
set contains a quasi-strict point satisfying those conditions. The isolated
corollary has a classical robustness analogue in Okada \cite{okada}; the
self-contained argument above supplies the particular continuous section
needed for the Mertens definition used here, rather than importing an
unstated equivalence of refinement terminology.
\end{remark}


\section{\texorpdfstring{A $2\times3$ obstruction: perfect points need not be stable singletons}{A 2 by 3 obstruction: perfect points need not be stable singletons}}
Consider the game
\begin{equation}\label{eq:2by3game}
\begin{array}{c|ccc}
 &L&M&R\\\hline
 U&(0,0)&(1,-1)&(-1,-1)\\
 D&(0,0)&(0,-1)&(0,-1).
\end{array}
\end{equation}
The column player strictly prefers $L$ to either other action against every
row mixture. At $L$ the row player is indifferent. Consequently,
\begin{equation}\label{eq:segment}
 E:=\NE(G)=\{z_t=(tU+(1-t)D,L):0\le t\le1\}.
\end{equation}

\begin{proposition}[All equilibria perfect, none singleton-stable]
\label{prop:obstruction}
For game \eqref{eq:2by3game}, every equilibrium $z_t$ is perfect (indeed
proper), but no singleton $\{z_t\}$ is Mertens $0$-stable. Every stable set
must equal the entire segment $E$.
\end{proposition}
\begin{proof}
For perfection, let $\epsilon\downarrow0$, put
\[
 p^\epsilon(U)=(1-2\epsilon)t+\epsilon,\qquad
 q^\epsilon=(1-2\epsilon,\epsilon,\epsilon).
\]
These profiles are completely mixed and tend to $z_t$. The row payoff
difference is $q^\epsilon(M)-q^\epsilon(R)=0$, and $L$ is always the column
player's best reply. Thus the target mixed strategies are best replies
against each perturbed opponent, proving perfection.

For the optional stronger claim of properness, replace the column profile
by $q^\epsilon=(1-2\epsilon^2,\epsilon^2,\epsilon^2)$, with
$0<\epsilon<1/2$, and retain the row profile above. The row actions tie.
Each inferior column action has probability $\epsilon^2$, whereas the
better action has probability $1-2\epsilon^2$, and
\[
 \epsilon^2\le\epsilon(1-2\epsilon^2).
\]
Thus these profiles are $\epsilon$-proper: whenever an action has strictly
lower payoff than another, its probability is at most $\epsilon$ times
the probability of the better action. Their limits are the claimed proper
equilibria.

Now impose row lower bounds $\alpha=(\delta/2,\delta/2)$. Consider the two
strictly positive column perturbations, ordered as $(L,M,R)$,
\[
 \beta^+=(\delta/4,\delta/2,\delta/4),\qquad
 \beta^-=(\delta/4,\delta/4,\delta/2).
\]
All players' lower bounds have common mass $\delta$. Since $L$ is strictly
best, the column player uniquely puts $M,R$ at their lower bounds and all
remaining mass on $L$. Under $\beta^+$ the row payoff difference is
$\delta/4>0$, giving the unique constrained equilibrium
\[
 p^+(U)=1-\delta/2,\qquad
 q^+=(1-3\delta/4,\delta/2,\delta/4).
\]
Under $\beta^-$ it is $-\delta/4<0$, giving
\[
 p^-(U)=\delta/2,\qquad
 q^-=(1-3\delta/4,\delta/4,\delta/2).
\]
The two equilibria converge respectively to $z_1$ and $z_0$.
Lemma~\ref{lem:necessary} and closedness imply that every stable set
must contain both $z_0$ and $z_1$: a compact set missing one of them has a
neighborhood avoided by its corresponding equilibrium ray for all small
$\delta$. Stable sets are connected and are subsets of $E$, so they must
contain the whole interval. No singleton can do so. Finally, existence of
Mertens stable sets in finite games, as used in \cite{gw}, implies that $E$
itself is the unique stable set. Only this last existence conclusion uses
that general existence theorem.
\end{proof}

At $0<t<1$, the equilibrium is quasi-strict with $I=\{U,D\}$ and $J=\{L\}$.
Choosing $D$ as row reference gives
\[
 M_A=\begin{pmatrix}0\\1\end{pmatrix},\qquad
 N_A=\begin{pmatrix}1&-1\\1&1\end{pmatrix}.
\]
Thus $\rank M_A=1$ whereas $\rank[M_A\ N_A]=2$.
Supported row indifference would require $\beta_M=\beta_R$, and a change in
the only supported column probability cannot compensate unequal trembles.
This is precisely the obstruction in Theorem~\ref{thm:rank}.

\begin{remark}[Not a counterexample to Econ-20]
The proposition disproves ``every perfect bimatrix equilibrium is a stable
singleton,'' even with ``proper'' in place of ``perfect.'' It does not
construct a rule $\Phi$ satisfying $B,D^*,I^*$ globally and selecting any
proper subinterval or singleton of $E$. Normal-form perfection or properness
at one game does not establish global invariance. Whether the exact axioms
force a particular selected set to expand to the stable set is a different
question. No assertion that this specific matrix defeats those axioms is made.
\end{remark}


\section{Two elementary comparisons that really satisfy the axioms}
Let a bimatrix game have row actions $\{1,\ldots,m\}$, column actions
$\{1,\ldots,n\}$ and payoff matrices $A,B$. For a column distribution
$\tau\in\Delta_n$ and $c>0$, add a column action $t$ whose payoffs are
\begin{equation}\label{eq:penalized}
 \bar A_{:,t}=A\tau,\qquad
 \bar B_{:,t}=B\tau-c\one.
\end{equation}
Multiple such actions may be added simultaneously.

\begin{lemma}[Penalized mixtures are an admissible $D^*$ extension]
\label{lem:dominated}
An extension formed by finitely many actions \eqref{eq:penalized}
is obtained by adding strongly irrelevant dominated strategies in the
source's exact sense.
\end{lemma}
\begin{proof}
Each $t$ is strictly dominated, for the column player, by the original
mixed strategy $\tau$, with the uniform gap $c$. Consequently no new
action can be used in a Nash equilibrium. Every original equilibrium
survives, since at it each new action is strictly worse than some original
best reply. Thus the two games have the same equilibria on the original
face.

Define $f_1(p)=p$ and
\[
 f_2(\bar q)=\bar q_{\rm old}+\sum_t\bar q_t\tau^t.
\]
This affine map is the identity on the old simplex. For any pair of
original row strategies $p,p'$, payoff comparisons against $\bar q$
are exactly the original comparisons against $f_2(\bar q)$, by the
first equality of \eqref{eq:penalized}. For any pair of original column
strategies, payoff comparisons against $p$ are unchanged. These statements
hold on the entire enlarged strategy space, so in particular on a
neighborhood of the original face. They verify every condition of
Definition 4.4, not merely dominance.
\end{proof}

\begin{lemma}[Duplicate plans and a hidden continuation]\label{lem:plans}
Fix a nonempty proper subset $J$ of original column actions and add actions
$t_j$, $j\in J$, as in \eqref{eq:penalized}. There is a finite perfect-recall
tree with the following properties. At its root the column player chooses
an action of $J$ directly or chooses $\mathrm{Out}$. After $\mathrm{Out}$
the same player chooses from $J^c\cup\{t_j:j\in J\}$. The row player
then acts at one information set, observing neither column decision.
The full normal form $H$ strongly embeds $\bar G$ by a duplicate-collapse
map $f$, and $f_1$ is the identity on row strategies.
\end{lemma}
\begin{proof}
Give the terminal node with actual row action $i$ and actual column action
$k$ its payoff $(\bar A_{ik},\bar B_{ik})$. Perfect recall holds because
the column player's second information set remembers the root choice,
and the row player has no earlier decision.

A pure column plan specifies a root action and an off-path continuation.
Every plan choosing a particular $j\in J$ directly is a duplicate of
$j$, regardless of its continuation. A plan choosing $\mathrm{Out}$
and then $k$ is identified with $k$. Collapsing these duplicates is an
affine surjection. Payoffs in $H$ are exactly those of $\bar G$ after
collapse. This is a strong embedding with no outsiders: the outsider face
is a point and the preference-equivalence condition holds globally.

All actual action pairs are distinct terminal outcomes. Since the row
player's information set is always reached, a mixed profile realization-
equivalent to a behavioral profile has the same collapsed row and column
marginals. Thus a $B$-witness in a lifted set projects into its specified
image in $\bar G$; the unreached continuation is not incorrectly treated
as an observed action.
\end{proof}


\section{A finite strategic exclusion certificate}
For a row distribution $p$, let $\BR_2(p)$ denote the column player's
pure best replies. The next theorem turns a linear payoff obstruction
into a valid global-axiom contradiction.

\begin{theorem}[Uniform indifference exclusion]\label{thm:certificate}
Let $K\subseteq\NE(A,B)$ be nonempty and compact. Suppose that there are
nonempty sets $I$ of row actions and $J$ of column actions, with $J^c$
nonempty, and a vector $w\in\R^m$ such that
\begin{align}
 &p_i>0\quad\text{for all }(p,q)\in K\text{ and }i\in I,\label{eq:used}\\
 &\BR_2(p)\subseteq J\quad\text{for all }(p,q)\in K,\label{eq:best}\\
 &\supp(w)\subseteq I,\qquad \one^{\mathsf T}w=0,\qquad
     w^{\mathsf T}A_{:,j}=0\quad(j\in J),\label{eq:annihilate}\\
 &w^{\mathsf T}A_{:,\ell}\ne0\quad\text{for some }\ell\in J^c.
 \label{eq:detect}
\end{align}
Then $K$ cannot be a selected set of any globally defined rule satisfying
$B,D^*,I^*$. The certificate uses at most $|J|$ new dominated actions,
a duplicate-plan embedding, and one application of $B$.
\end{theorem}
\begin{proof}
Put $v(p)=\max_j p^{\mathsf T}B_{:,j}$. By \eqref{eq:best}, compactness,
continuity, and finiteness of the action sets,
\[
 \gamma=\min_{(p,q)\in K}\min_{k\in J^c}
             \{v(p)-p^{\mathsf T}B_{:,k}\}>0.
\]
Let $L=\max_{i,j,k}|B_{ij}-B_{ik}|$. Then $L\ge\gamma>0$.
Choose
\[
 \epsilon=\frac{\gamma}{4L},\qquad c=\frac{\gamma}{4},\qquad
 \tau^j=(1-\epsilon)e_j+\epsilon e_\ell\quad(j\in J).
\]
Add a penalized mixture $t_j$ for each $\tau^j$. Lemma
\ref{lem:dominated} applies.

For each $(p,q)\in K$, choose $j\in\BR_2(p)\subseteq J$. The payoff
from its associated continuation action satisfies
\begin{align*}
 p^{\mathsf T}\bar B_{:,t_j}
 &= (1-\epsilon)v(p)+\epsilon p^{\mathsf T}B_{:,\ell}-c\\
 &\ge v(p)-\epsilon L-c
 =v(p)-\gamma/2
 \ge p^{\mathsf T}B_{:,k}+\gamma/2\quad(k\in J^c).
\end{align*}
Thus every original action in $J^c$ is strictly worse than some new action
at the continuation information set. The inequality persists against row
profiles sufficiently close to the projection of $K$.

Assume for contradiction that $K\in\Phi(G)$. Axiom $D^*$ makes $K$,
padded by zero probabilities, a selected set of $\bar G$. Apply
Lemma \ref{lem:plans} and the existence-of-lifts direction of $I^*$ to
obtain $K_H\in\Phi(H)$ with $f(K_H)=K$. Axiom $B$ supplies a profile
in $K_H$ realization-equivalent to a quasi-perfect behavioral equilibrium
$b$ of the hidden-continuation tree. Its collapsed image lies in $K$.
Let $b^r\to b$ be the completely mixed behavioral sequence in the
definition of quasi-perfection.

The continuation's limiting strategy puts total probability one on the
new actions: any positive limiting probability on an old $k\in J^c$
would contradict its strict suboptimality against the approximating row
strategies. Denote its probabilities along $b^r$ by $\zeta_k^r$, and
let $\theta_r>0$ be the root probability of $\mathrm{Out}$. We have
\[
 \sum_{j\in J}\zeta_{t_j}^r\longrightarrow1,
 \qquad \zeta_k^r\longrightarrow0\quad(k\in J^c).
\]
Writing $a=w^{\mathsf T}A_{:,\ell}\ne0$, equation
\eqref{eq:annihilate} gives the same value for every new action:
\[
 w^{\mathsf T}\bar A_{:,t_j}=\epsilon a.
\]
All root-direct actions have weighted row payoff zero. The weighted
row payoff against the full induced column distribution $\bar q^r$ is
therefore
\begin{equation}\label{eq:contradiction}
 w^{\mathsf T}\bar A\bar q^r
 =\theta_r\left(
 \epsilon a\sum_{j\in J}\zeta_{t_j}^r
 +\sum_{k\in J^c}\zeta_k^r w^{\mathsf T}A_{:,k}\right).
\end{equation}
The bracket converges to $\epsilon a\ne0$; since $\theta_r>0$, the
expression is nonzero for all sufficiently large $r$.

On the other hand, the target row strategy uses every action in $I$.
Quasi-perfection at its sole information set makes each such action a
best reply to the opponent's behavioral strategy $b_2^r$ for every $r$.
Their payoffs are equal. Since $w$ is supported on $I$ and its coordinates
sum to zero, their weighted payoff must be zero. This contradicts
\eqref{eq:contradiction}. No globally compatible rule can select $K$.
\end{proof}

\begin{remark}[What is and is not being assumed]
The new actions are strictly dominated in the original normal-form sense,
yet optimal after an off-path restriction to continuation choices. There
is no contradiction: the dominating original mixture uses root-direct
actions that are unavailable once $\mathrm{Out}$ has been chosen.
The proof retains the distinction between a simultaneous game's perfect
equilibrium and quasi-perfection in all its allowed representations.
It does not assume that equilibrium selection is continuous in payoffs.
\end{remark}


\section{Selected sets consisting entirely of quasi-strict equilibria}
An equilibrium $(p,q)$ is \emph{quasi-strict} when
$\BR_1(q)=\supp(p)$ and $\BR_2(p)=\supp(q)$: unused actions are
strictly inferior. This is not the definition of quasi-perfect equilibrium.

\begin{lemma}[Supports are constant on a connected quasi-strict set]
\label{lem:supports}
If a connected set $K\subseteq\NE(G)$ consists entirely of quasi-strict
equilibria, the two supports $I=\supp(p)$ and $J=\supp(q)$ do not depend
on $(p,q)\in K$.
\end{lemma}
\begin{proof}
At any quasi-strict equilibrium, positive coordinates remain positive in
a sufficiently small neighborhood. Every unused action remains strictly
worse than an originally supported action in that neighborhood, so it
cannot be used at a neighboring Nash equilibrium. Hence the support map
is locally constant on $K$. It takes values in a finite discrete set;
connectedness forces it to be constant.
\end{proof}

Fix reference actions $i_0\in I$ and $j_0\in J$, and set
\[
 C_{i,:}=A_{i,:}-A_{i_0,:}\quad(i\in I\setminus\{i_0\}),\qquad
 D_{j,:}=(B_{:,j}-B_{:,j_0})^{\mathsf T}\quad(j\in J\setminus\{j_0\}).
\]
Allow empty difference matrices. Define
\begin{equation}\label{eq:matrices}
 \begin{aligned}
 M_A&=\begin{pmatrix}C_{:,J}\\\one_J^{\mathsf T}\end{pmatrix},&
 N_A&=\begin{pmatrix}C_{:,J^c}\\\one_{J^c}^{\mathsf T}\end{pmatrix},\\
 M_B&=\begin{pmatrix}D_{:,I}\\\one_I^{\mathsf T}\end{pmatrix},&
 N_B&=\begin{pmatrix}D_{:,I^c}\\\one_{I^c}^{\mathsf T}\end{pmatrix}.
 \end{aligned}
\end{equation}

\begin{theorem}[Quasi-strict selected sets are stable]\label{thm:qs}
Let $\Phi$ satisfy the global axioms and let $K\in\Phi(G)$ for a finite
bimatrix game. If every point of $K$ is quasi-strict, then
\[
 \im N_A\subseteq\im M_A,\qquad \im N_B\subseteq\im M_B.
\]
Moreover, $K$ itself is Mertens $0$-stable, and every $z\in K$ has a
stable singleton. In particular the target containment conclusion holds
for this selected set, without any isolation or genericity assumption.
\end{theorem}
\begin{proof}
Supports are constant by Lemma \ref{lem:supports}. Suppose first that
$\im N_A\not\subseteq\im M_A$. A separating row vector can be written
$(h^{\mathsf T},\lambda)$, with
\[
 (h^{\mathsf T},\lambda)M_A=0,\qquad
 (h^{\mathsf T},\lambda)N_A\ne0.
\]
Choose any $(p^*,q^*)\in K$. Since $C_{:,J}q_J^*=0$ and
$\one^{\mathsf T}q_J^*=1$, multiplication by $q_J^*$ shows
$\lambda=0$. Consequently
\[
 h^{\mathsf T}C_{:,J}=0,\qquad h^{\mathsf T}C_{:,\ell}\ne0
 \quad\text{for some }\ell\in J^c.
\]
Define $w_i=h_i$ for $i\in I\setminus\{i_0\}$,
$w_{i_0}=-\sum_{i\ne i_0}h_i$, and $w_i=0$ outside $I$.
Then $\one^{\mathsf T}w=0$ and $w^{\mathsf T}A=h^{\mathsf T}C$.
All hypotheses of Theorem \ref{thm:certificate} hold, since each point
of $K$ uses $I$ and has column best replies exactly $J$. This contradicts
selection of $K$. The first inclusion follows. Exchanging the two players
proves the second. Empty complements cause no difficulty; their inclusions
are automatic.

Choose matrices $R_A,R_B$ solving $M_AR_A=N_A$, $M_BR_B=N_B$.
For $z=(p^*,q^*)\in K$ and row/column lower bounds $\alpha,\beta$, put
\begin{equation}\label{eq:uniformsection}
 \begin{aligned}
 q_{J^c}&=\beta_{J^c},&q_J&=q_J^*-R_A\beta_{J^c},\\
 p_{I^c}&=\alpha_{I^c},&p_I&=p_I^*-R_B\alpha_{I^c}.
 \end{aligned}
\end{equation}
The probability-sum rows of the identities make $p,q$ distributions;
the other rows give $Cq=0$, $Dp=0$. Thus supported actions stay tied.

All coordinates in $p_I^*,q_J^*$ have a common positive lower bound as
$z$ ranges over compact $K$. Likewise the unused-action payoff gaps have
a common positive lower bound whenever such actions exist. The changes
in \eqref{eq:uniformsection} are bounded uniformly by a constant times
the perturbation mass. Hence, for a common sufficiently small $\rho>0$,
all supported coordinates exceed their required lower bounds and all
unused actions remain strictly inferior, uniformly over $K$.
Inferior actions receive exactly their lower bounds; all remaining
probability is assigned to tied best replies. Therefore
\eqref{eq:uniformsection} is a jointly continuous family of constrained
equilibria. Lemma \ref{lem:family} proves both conclusions about stability.
\end{proof}

\begin{corollary}[A necessary feature of any counterexample]
\label{cor:boundary}
If a selected set $K$ in a bimatrix game contains no stable subset, then
$K$ contains at least one non-quasi-strict equilibrium. Equivalently,
at some point of $K$, an unused action is also a best reply.
\end{corollary}
\begin{proof}
Otherwise Theorem \ref{thm:qs} makes $K$ itself stable.
\end{proof}


\section{Completing the earlier \texorpdfstring{$2\times3$}{2 x 3} example at the set level}
Consider
\begin{equation}\label{eq:example}
 \begin{array}{c|ccc}
  &L&M&R\\\hline
  U&(0,0)&(1,-1)&(-1,-1)\\
  D&(0,0)&(0,-1)&(0,-1).
 \end{array}
\end{equation}
The Nash set is $E=\{z_t=(tU+(1-t)D,L):0\le t\le1\}$.
Column $L$ is strictly dominant, and row actions tie at $L$.

\begin{theorem}[The axioms force the entire segment]\label{thm:segment}
For every globally compatible rule $\Phi$, game \eqref{eq:example}
satisfies $\Phi(G)=\{E\}$. Moreover, $\Stab(G)=\{E\}$.
Thus this game is not a counterexample to the two-player conjecture,
even though no individual equilibrium has a stable singleton.
\end{theorem}
\begin{proof}
First add a new column action $T^+$ with payoffs
\[
 \bar G(U,T^+)=(1/2,-3/4),\qquad
 \bar G(D,T^+)=(0,-3/4).
\]
It is the penalized mixture $\tau^+=(L+M)/2$, with penalty $1/4$.
Lemma \ref{lem:dominated} makes this an exact $D^*$ extension.
Use the tree with root choices $L,\mathrm{Out}$ and continuation
choices $M,R,T^+$, followed by the uninformed row decision. Its full
normal form has six column plans, three of them duplicates of $L$.
Lemma \ref{lem:plans} and $I^*$ justify the collapse.

At the continuation information set, $T^+$ gives the column player
$-3/4$, strictly exceeding the payoff $-1$ of either alternative,
against every row strategy. A quasi-perfect equilibrium must therefore
choose $T^+$ there. At the root it chooses $L$, which gives $0$.
In a completely mixed behavioral sequence witnessing quasi-perfection,
write $\theta_r>0$ for the probability of $\mathrm{Out}$ and
$\zeta^r$ for the continuation distribution. Then
\[
 u_1(U,b_2^r)-u_1(D,b_2^r)
 =\theta_r(\zeta_M^r-\zeta_R^r+\tfrac12\zeta_{T^+}^r)>0
\]
for all sufficiently large $r$, since $\zeta_{T^+}^r\to1$.
Every quasi-perfect outcome of this tree is therefore $(U,L)$.
Lifting an arbitrary $K\in\Phi(G)$ through $D^*$ and the duplicate
embedding, then applying $B$, forces $z_1\in K$.

Now instead add the penalized mixture $\tau^-=(L+R)/2$ with the same
penalty. Its row payoffs are $-1/2,0$ and its column payoff is again
$-3/4$. The identical continuation argument makes the row payoff difference
strictly negative along every quasi-perfect witness, so it forces
$z_0=(D,L)$ into the \emph{same original selected set} $K$.
A nonempty connected subset of the segment $E$ containing its two
endpoints is $E$. Thus every selected set is $E$; nonemptiness of the
solution family gives $\Phi(G)=\{E\}$.

Apply this proved selection statement to the known all-game solution
concept $\Phi=\Stab$. Its existence and satisfaction of the axioms are
independent literature results cited in Section 1. It follows that
$\Stab(G)=\{E\}$. This application does not invoke the unproved
nongeneric containment converse.
\end{proof}

\begin{remark}[Direct perturbation check]
For row lower bounds $(\delta/2,\delta/2)$, column lower bounds
$(\delta/4,\delta/2,\delta/4)$ give the unique constrained equilibrium
with row $U$-probability $1-\delta/2$. Interchanging the last two column
bounds gives the unique constrained equilibrium with row $U$-probability
$\delta/2$. These two limits independently explain why a singleton
cannot be robust to all tremble directions. They do not themselves
establish that a globally selected set contains both limits; the two
admissible extensions in the theorem supply that missing implication.
\end{remark}


\section{An unbounded-size subclass: a strictly dominant column}
Let the row player have two actions $U,D$ and the column player have
arbitrary finite actions $0,1,\ldots,n-1$. Suppose
\begin{equation}\label{eq:dominant}
 B_{i0}>B_{ij}\quad(i\in\{U,D\},\ j\ne0).
\end{equation}
Define $d_j=A_{Uj}-A_{Dj}$.

\begin{theorem}[Containment for every dominant-column $2\times n$ game]
\label{thm:dominant}
Every game satisfying \eqref{eq:dominant} satisfies the target
\eqref{eq:target}. In the special case
\begin{equation}\label{eq:oppositesigns}
 d_0=0,\qquad \min_{j\ne0}d_j<0<\max_{j\ne0}d_j,
\end{equation}
every compatible rule selects exactly the full Nash segment
$E=\{(tU+(1-t)D,0):0\le t\le1\}$, and this segment is the unique
Mertens-stable set.
\end{theorem}
\begin{proof}
Column $0$ is the unique column best reply to every row distribution.
If $d_0>0$ or $d_0<0$, there is a unique strict Nash equilibrium.
For all sufficiently small lower bounds, keep the dominant column
at all its feasible residual mass and clamp the strict row action
to its feasible endpoint. The strict payoff gap persists, yielding a
continuous equilibrium section. Its singleton is stable by
Lemma \ref{lem:family}. This settles $d_0\ne0$.

Now let $d_0=0$, so the Nash set is $E$. If all $d_j=0$, every target
$t\in[0,1]$ has the section
\[
 p_U=\alpha_U+(1-\delta)t,\quad
 p_D=\alpha_D+(1-\delta)(1-t),\quad
 q_j=\beta_j\ (j\ne0),\quad q_0=1-\sum_{j\ne0}\beta_j.
\]
Both row actions always tie, and column $0$ remains strictly best.
Every singleton is stable, so every nonempty selected set contains one.

If all $d_j\ge0$ and at least one is positive, $U$ is the unique best
reply against every completely mixed column strategy. Axiom $B$, applied
to the simultaneous-move representation, therefore places $(U,0)$
in each selected set. The formula $p_U=1-\alpha_D$, together with the
column formula above, is a continuous equilibrium section because
$\sum_jd_jq_j\ge0$ for every nonnegative lower-bound vector. Its zero
fiber is $(U,0)$, proving the requisite stability. The case with all
$d_j\le0$ and at least one negative is symmetric.

It remains to treat \eqref{eq:oppositesigns}. Put
\[
 \gamma=\min_{i,j\ne0}(B_{i0}-B_{ij})>0,\qquad
 M=\max_{i,j\ne0}(B_{i0}-B_{ij})\ge\gamma,
 \qquad \epsilon=\frac{\gamma}{4M},\quad c=\frac\gamma4.
\]
Choose $h$ with $d_h>0$, and add the penalized column mixture
$\tau=(1-\epsilon)e_0+\epsilon e_h$. For every row action $i$
and original nonzero column $j$,
\begin{align*}
 \bar B_{it}-B_{ij}
 &= B_{i0}-B_{ij}-\epsilon(B_{i0}-B_{ih})-c\\
 &\ge\gamma-\epsilon M-c=\gamma/2>0.
\end{align*}
Column $0$ strictly beats $t$, and $t$ strictly beats every original
nonzero column. Its row payoff difference is $\epsilon d_h>0$.
The root choice $0/\mathrm{Out}$ followed by all nonzero original
columns and $t$ is therefore exactly the forcing construction of
Theorem \ref{thm:segment}. It forces $(U,0)$ into every selected $K$.
Repeating with a column $h$ satisfying $d_h<0$ forces $(D,0)$ into the
same $K$. Connectedness gives $K=E$. Finally apply the proved conclusion
to the known nonempty stable-set solution concept to get
$\Stab(G)=\{E\}$.

The cases $d_0\ne0$, all differences zero, one weak sign, and both
strict signs exhaust all possibilities. The proof permits any finite
$n$; $n=1$ falls within the first or all-zero case.
\end{proof}


\section{A dominant column and a balanced family of payoff vectors}
Let the row actions be $\{1,\ldots,m\}$ and the column actions be
$\{0,1,\ldots,n-1\}$, with $n\ge2$. Assume column $0$ is strictly
dominant:
\begin{equation}\label{eq:dominance}
 B_{i0}>B_{ij}\qquad\text{for all }i\text{ and }j\ne0.
\end{equation}
Set
\[
 I=\arg\max_i A_{i0},\qquad s=|I|,\qquad
 E=\Delta(I)\times\{e_0\}.
\]
Here $\Delta(I)$ is the face of the row simplex supported on $I$.
By strict dominance and the row best-response condition,
\begin{equation}\label{eq:ne}
 \NE(G)=E.
\end{equation}
Choose any finite list of original column mixtures
$\tau^1,\ldots,\tau^k\in\Delta_n$ and define
\[
 C_{ij}=(A\tau^j)_i\quad(1\le i\le m,\ 1\le j\le k),
 \qquad C_I=(C_{ij})_{i\in I,\,j\le k}.
\]
The following are explicit sufficient hypotheses, not consequences of
the three axioms:
\begin{align}
 &\text{There exist }z\in\Delta_k,\ z_j>0\text{ for every }j,
   \text{ and }a\in\R\text{ such that }C_Iz=a\one_I;
   \label{eq:balance}\\
 &w\in\R^I,\quad\one_I^Tw=0,\quad C_I^Tw=0
   \quad\Longrightarrow\quad w=0.\label{eq:span}
\end{align}
We call \eqref{eq:balance} balance and \eqref{eq:span} spanning.
Equivalently, spanning says
\[
 \rank\begin{pmatrix}C_I^T\\\one_I^T\end{pmatrix}=s.
\]
Balance supplies a strictly positive mixture of probe columns that
makes all rows in $I$ indifferent. Spanning says that these probe
payoffs distinguish row distributions on that face.


\section{A calibrated, admissible extension}
Fix an arbitrary target $r\in\Delta(I)$, allowing zero coordinates.
Write
\[
 a_j(r)=\sum_{h\in I}r_hC_{hj},\qquad
 \gamma=\min_{i,\,\ell\ne0}(B_{i0}-B_{i\ell})>0,
\]
and choose
\[
 D_r=\max\{1,\max_{i,j}|a_j(r)-C_{ij}|\},\qquad
 c=\gamma/2,\qquad\kappa=\frac{\gamma}{8D_r}>0.
\]
Add column actions $t_1,\ldots,t_k$, retaining every original payoff,
and set
\begin{equation}\label{eq:new}
 \bar A_{i,t_j}=C_{ij},\qquad
 \bar B_{i,t_j}=B_{i0}-c+\kappa\bigl(a_j(r)-C_{ij}\bigr).
\end{equation}

\begin{lemma}[Exact admissibility]\label{lem:admissible}
The extension \eqref{eq:new} satisfies the source's $D^*$ comparison.
Moreover, every new action lies strictly below column $0$ and strictly
above every original nonzero column in column-player payoff, uniformly
in the row player's strategy.
\end{lemma}
\begin{proof}
For every pure row $i$ and probe $t_j$,
\[
 \frac{3\gamma}{8}\le B_{i0}-\bar B_{i,t_j}
                       \le\frac{5\gamma}{8}.
\]
For every original $\ell\ne0$,
\[
 \bar B_{i,t_j}-B_{i\ell}
 =(B_{i0}-B_{i\ell})-(B_{i0}-\bar B_{i,t_j})
 \ge\frac{3\gamma}{8}.
\]
The same bounds hold against mixed rows by linearity. Thus every new
action is strictly dominated by the original pure action $0$; column
$0$ remains strictly dominant in the enlarged game. Its Nash set is
exactly the original face $E$, with zero probabilities on new actions.
This verifies Definitions 4.1--4.2 of the source.

For Definition 4.4 use the affine dictionary
\[
 f_1(p)=p,\qquad
 f_2(\bar q)=\bar q_{\rm old}+\sum_{j=1}^k\bar q_{t_j}\tau^j.
\]
It is the identity on the original face. For original row strategies
$p,p'$, the difference in their enlarged-game payoffs against $\bar q$
is exactly $(p-p')^TAf_2(\bar q)$, because the new row-payoff columns
are $A\tau^j$. For any two original column strategies, their payoff
comparison is unchanged and evaluated against $f_1(p)=p$.
The required preference equivalences therefore hold everywhere, not
merely in a neighborhood. Together with domination and equality of
Nash sets, these facts establish the exact $D^*$ comparison.
\end{proof}


\section{Full normal form, lifting, and quasi-perfection}
Represent $\bar G$ by a finite perfect-recall tree. The column player
first chooses $0$ or Out. After Out, that same player chooses one of
\[
 \{1,\ldots,n-1,t_1,\ldots,t_k\}.
\]
The row player then moves at a single information set, observing neither
column move, with all $m$ row actions available. Payoffs are those of the
resulting row and column in $\bar G$.

Let $H$ be the \emph{full} normal form. A column pure plan specifies
both a root action and a continuation; hence there are
$n-1+k$ duplicate plans with root action $0$. Collapsing those plans
to $0$, and mapping each Out plan to its continuation action, gives an
affine surjection $g$ from the strategy space of $H$ onto that of
$\bar G$. It preserves both payoff functions everywhere. With no
outsiders, this is a strong embedding under Definition 5.2.

Consequently, if $K\in\Phi(G)$, Lemma \ref{lem:admissible} and $D^*$
put the padded $K$ in $\Phi(\bar G)$. Axiom $I^*$ supplies
$K_H\in\Phi(H)$ with $g(K_H)=K$. Axiom $B$ supplies a profile in
$K_H$ equivalent to a quasi-perfect equilibrium $b$ of this tree.
The row mixture $p$ of $b$ is therefore a row mixture appearing in
$K$. The root action is $0$: it is strictly better than every possible
continuation. In particular $p\in\Delta(I)$.

\begin{lemma}[The induced continuation inequalities]\label{lem:continuation}
For that quasi-perfect equilibrium there is $y\in\Delta_k$ such that
\begin{align}
 &p^TC_Iy\ge r^TC_Iy,\label{eq:rowbr}\\
 &\supp(y)\subseteq
       \arg\min_j (p-r)^TC_{I,j}.\label{eq:colbr}
\end{align}
\end{lemma}
\begin{proof}
Let $b^\nu\to b$ be a completely mixed behavioral sequence witnessing
quasi-perfection. All original nonzero continuation actions are uniformly
inferior to the new actions, by Lemma \ref{lem:admissible}. None can have
positive probability in the target continuation. Thus the target
continuation is a distribution $y$ over $t_1,\ldots,t_k$.

At the continuation node, every target action with positive probability
is optimal against the row component of $b^\nu$. Passing to the limit
and deleting the common term $p^TB_{:,0}-c$ from \eqref{eq:new} gives
\eqref{eq:colbr}.

For completeness the row-player argument must account for the fact that
Out is unreached at $b$. Let $\theta_\nu>0$ be its probability in
$b^\nu$, and let $\zeta^\nu$ be the conditional continuation distribution.
Its mass on original nonzero columns tends to zero, while its probe
probabilities tend to $y$. For $i\in\supp(p)$ and $h\in I$,
quasi-perfection requires the payoff difference of row $i$ over row $h$
to be nonnegative against $b_2^\nu$. Both rows have the same payoff
against root action $0$. Dividing the inequality by $\theta_\nu>0$
and taking the limit yields
\[
 (C_Iy)_i\ge(C_Iy)_h.
\]
Every row used by $p$ is therefore a best reply to $y$ among $I$.
Comparing with $r\in\Delta(I)$ gives \eqref{eq:rowbr}. No comparison
between two vanishing error rates has been assumed: the continuation
probabilities are conditioned on the single positive factor
$\theta_\nu$, which cancels exactly.
\end{proof}


\section{Payoff-fiber forcing and full-face selection}
\begin{theorem}[Payoff-fiber forcing]\label{thm:fiber}
Assume \eqref{eq:dominance} and \eqref{eq:balance}, but not necessarily
spanning. For every globally $B,D^*,I^*$-compatible $\Phi$, every
$K\in\Phi(G)$, and every $r\in\Delta(I)$, there exists
$(p,e_0)\in K$ with
\begin{equation}\label{eq:fiber}
 C_I^Tp=C_I^Tr.
\end{equation}
Equivalently, the row projection of $K$ maps onto the entire polytope
$C_I^T\Delta(I)$ under $p\mapsto C_I^Tp$.
\end{theorem}
\begin{proof}
Use the calibrated extension for the chosen $r$ and the lifted
quasi-perfect equilibrium from the previous section. Set
$d_j=(p-r)^TC_{I,j}$. By balance,
\[
 \sum_jz_jd_j=(p-r)^TC_Iz
 =a\one_I^T(p-r)=0.
\]
Thus $\min_jd_j\le0$. By \eqref{eq:colbr}, $y$ is supported on
minimizers of $d$, whereas \eqref{eq:rowbr} says
\[
 0\le\sum_jy_jd_j=\min_jd_j\le0.
\]
It follows that all $d_j\ge0$. Their average under the strictly positive
vector $z$ is zero, so every $d_j=0$. This proves \eqref{eq:fiber}.
The point $p$ lies in the original selected $K$ by the exact lifting
argument. Since $r$ was arbitrary, the image equality follows.
\end{proof}

\begin{theorem}[Exact selection of a balanced-spanning equilibrium face]
\label{thm:face}
Assume \eqref{eq:dominance}, \eqref{eq:balance}, and \eqref{eq:span}.
Then every globally compatible solution concept satisfies
\[
 \boxed{\Phi(G)=\{E\},\qquad E=\Delta(I)\times\{e_0\}.}
\]
Moreover,
\[
 \boxed{\Stab(G)=\{E\}.}
\]
In particular, the containment conclusion of WS-107793 holds for this
class, with no bound on the numbers of players' actions or on $\dim E$.
\end{theorem}
\begin{proof}
Fix $K\in\Phi(G)$ and $r\in\Delta(I)$. Theorem \ref{thm:fiber}
provides $(p,e_0)\in K$ with $C_I^T(p-r)=0$. Since
$\one_I^T(p-r)=0$, spanning gives $p=r$. Thus $(r,e_0)\in K$ for
\emph{every} $r\in\Delta(I)$. As $K\subseteq\NE(G)=E$, $K=E$.
Nonemptiness of $\Phi(G)$ gives the claimed equality of families.

Apply that equality to the source's stability solution concept itself.
The independent source inputs specified in Section 1 say that it exists
and satisfies the three axioms on all games. Therefore it selects
exactly $E$ in this game. This identifies the genuine Mertens-stable
set, not merely a topologically robust substitute. The nongeneric
converse (T) has not been assumed at any stage.
\end{proof}

\begin{corollary}[One finite comparison chain excludes every proper subset]
Under the hypotheses of Theorem \ref{thm:face}, every nonempty proper
subset $K\subsetneq E$ is excluded by the particular extension calibrated
to any $r$ with $(r,e_0)\notin K$. The chain uses $k$ added dominated
column actions, one duplicate-plan embedding, and one application of $B$.
\end{corollary}
\begin{proof}
Were $K$ selected, the preceding argument would force $(r,e_0)\in K$.
The construction is finite and uses exactly those comparisons.
\end{proof}

\begin{remark}[Why this is more than endpoint forcing]
The proof does not infer that a connected set containing the vertices
of a simplex contains its interior; that inference is false for dimension
at least two. Instead it calibrates a different allowed test to each
individual $r$. Indeed, connectedness of $K$ is not needed in this
particular forcing argument, although it remains part of the original
solution-concept definition.
\end{remark}


\section{An explicit family of every equilibrium dimension}
For each integer $d\ge2$, define a $d\times(d+1)$ game $G_d$ with
columns $0,1,\ldots,d$ by
\begin{equation}\label{eq:canonical}
 A_{i0}=0,\quad A_{ij}=\mathbf1_{\{i=j\}}\ (j\ge1),\qquad
 B_{i0}=0,\quad B_{ij}=-1\ (j\ge1).
\end{equation}
Then $I=\{1,\ldots,d\}$, $\NE(G_d)=\Delta_d\times\{e_0\}$.
Taking $\tau^j=e_j$ gives $C_I=\operatorname{Id}_d$ and
$z=d^{-1}\one$. Both balance and spanning hold. Hence
\[
 \Phi(G_d)=\Stab(G_d)=\{\Delta_d\times\{e_0\}\}.
\]
The equilibrium face has dimension $d-1$. None of its points is strict
or isolated. When $d\ge2$, no singleton is stable, since the only stable
set is the entire face.

For $d=3$ the original game is
\[
\begin{array}{c|cccc}
 &0&1&2&3\\\hline
1&(0,0)&(1,-1)&(0,-1)&(0,-1)\\
2&(0,0)&(0,-1)&(1,-1)&(0,-1)\\
3&(0,0)&(0,-1)&(0,-1)&(1,-1)
\end{array}.
\]
Fix any $r\in\Delta_3$. It suffices to add three column actions with
\[
 \bar A_{i,t_j}=\mathbf1_{\{i=j\}},\qquad
 \bar B_{i,t_j}=-\frac12+\frac18
                       \left(r_j-\mathbf1_{\{i=j\}}\right).
\]
Their column payoffs lie between $-5/8$ and $-3/8$, hence strictly
between the old nonzero columns and column $0$.
The limiting continuation satisfies
\[
 p\in\BR_1(y)\quad\text{for payoff }p^Ty,
 \qquad
 y\in\arg\max_{v\in\Delta_3}(r-p)^Tv.
\]
Equivalently it is a zero-sum equilibrium after subtracting $r^Ty$
from the row player's payoff, which is independent of his action.
The elementary calculation in Theorem \ref{thm:fiber} forces $p=r$.
Thus every interior point, not merely each vertex, lies in every
selected set.

\paragraph{A concrete numerical target.}
For $r=(1/2,1/3,1/6)$ the column-player payoff vectors of the probes are
\[
 \begin{aligned}
 \bar B_{:,t_1}&=(-9/16,-7/16,-7/16)^T,\\
 \bar B_{:,t_2}&=(-11/24,-7/12,-11/24)^T,\\
 \bar B_{:,t_3}&=(-23/48,-23/48,-29/48)^T.
 \end{aligned}
\]
The target continuation is uniform in this interior example, and its
row mixture is precisely $r$.


\section{Remaining gap, checks, and limitations}
The conclusions above do not prove the full implication
\eqref{eq:target} for arbitrary bimatrix targets. The rank theorem
concerns a \emph{singleton} at a quasi-strict equilibrium. The finite
exclusion certificate handles selected sets that retain a common
positive support and a uniform strictly inferior outside column.
The calibrated action construction requires a common strictly dominant
column and the stated balance--spanning assumptions. In an arbitrary
selected equilibrium set, the supports and best-reply faces may change;
unused actions can tie at the boundary. The strict gap used to place
new probes below a root action yet above old continuations can vanish.
Neither connectedness alone nor simple robustness to perturbations
replaces the rational \v Cech essentiality and connexity in
Definition 2.4 of~\cite{gw}. We have neither glued the available local
certificates into an unrestricted witness nor constructed a solution
concept satisfying $B,D^*,I^*$ globally that violates containment.

\paragraph{Scope of checking.}
All results are theoretical implications whose proofs appear above.
The three underlying drafts report finite exact-arithmetic checks of
selected payoff identities, ranks, and perturbation formulas. Those
computer calculations are not reproduced or submitted here, and this
manuscript does not claim that finite sampling checks any universal
statement. The proofs have been consolidated and the \LaTeX\ compilation
checked; no proof-assistant formalization or independent referee
validation is claimed. All arithmetic or game-theoretic hypotheses
needed for each result are displayed in its theorem statement.

\paragraph{Model and human provenance.}
The original three supplied manuscripts identify OpenAI GPT-6 Astra Pro
as an LLM used for proof search and preparation, with exact backend
version and thinking settings not available. This merging and editorial
review used OpenAI GPT-6. The submitter selected the problem, prompted
the research and requested consolidation. No independent human proof
review is claimed. Before submitting, the separate We-Solve form's
model-use, human roles, rights and public-disclosure declarations must
be completed accurately by the submitter; this paper does not supply
unknown personal identities or guess model settings.

\begin{thebibliography}{9}
\bibitem{gw}
Srihari Govindan and Robert Wilson.
\emph{Axiomatic Equilibrium Selection in Generic Extensive-Form Games}.
arXiv:2504.16908v2, 26 August 2026.
Relevant locators: Definition 2.1 (quasi-perfection), Definition 2.4
(stability), Definitions 4.4 and 5.2, Axioms 3.1, 4.7 and 5.4,
Theorem 6.1, Section 10 and Appendices A--B.
\url{https://arxiv.org/html/2504.16908v2}.

\bibitem{gm}
Srihari Govindan and Jean-Fran\c{c}ois Mertens.
\emph{An equivalent definition of stable equilibria}.
\emph{International Journal of Game Theory} \textbf{32} (2004), 339--357.
\url{https://doi.org/10.1007/s001820400165}.

\bibitem{okada}
A. Okada.
\emph{Strictly perfect equilibrium points of bimatrix games}.
\emph{International Journal of Game Theory} \textbf{13} (1984), 145--153.
\url{https://doi.org/10.1007/BF01769812}.

\bibitem{ws}
We-Solve research catalogue.
\emph{Axiomatic containment of Mertens-stable sets without genericity}.
Submission handle \texttt{Econ-20}; public research-record ID
\texttt{prb\_34ecf5fa262f4a85a0fbe4af9837c46b}.
\url{https://we-solve.github.io/problems/axiomatic-containment-of-mertens-stable-sets-without-genericity/}.
\end{thebibliography}
\end{document}
